% Nouveau EGA V (cote 25): archive cote 25, site page 48 (PDF page 49).
% typescript. Draft compared: none.
% AI-transcribed by Claude Opus 5.5 (claude-opus-5-5, adaptive thinking, effort high), Message Batches API, batch ega25b (msgbatch_01A7n5YuXubUPNK1zsepi8RG), 2026-11-13; not yet checked by eye.
\scanpage{25/48}
\mspage{\unsure{45|46}}

\deleted{pour toute composante irréductible $X'_i$ de $X'$ on a}
\inserted{grâce au fait que $\underline{X'}$ sing est de dim $r-1$, où $r=\dim \check{P}$ !}

que $\phi_{\hand{$\zeta$}}$ soit $X_{k(\hand{$\zeta$})}$-régulière et $Y_{\hand{$\zeta$}}$ satisfasse la condition $R_k$ géométri-%

que. Alors:

a) \over{B}{P}our que $Z_{k-1}$ soit de codimension $\geq 2$ dans $\overset{\hand{$\vee$}}{P}$ , il \xout{f et} s que
\inserted{$X'_i$ de $X'$ de dim $\leq k$ soit \uline{réglée} pour $f$.} \deleted{\unsure{et $X'$ qui est non réglée rel. à $f$}}
toute composante irréductible \deleted{de $X_i$ soit de codimension $\geq k+1$ et c'est aus-%
si nécessaire si codim$(f(\underline{Y}'_1{}^{\mathrm{sing}}),P) = 1$ (``cas non réglé'').}

b) Pour que $Z_k$ soit de codimension $\geq 2$ dans $\check{P}$ , il \xout{f et} s que
\deleted{\inserted{$X'_i$ \unsure{qui est non réglée rel. à $f$}}}
toute composante irréductible \deleted{de $X$ soit de dimension $\geq k+1$, et que}
\inserted{$X_i$ de dim $\leq k+1$ de $X$ soit formée de pts lisses de $X$ et soit \uline{réglée} pour $f$.}
\deleted{l'ensemble singulier $T = X-X'$ n'ait aucun point isolé $x$ avec}
\inserted{\unsure{i.e. toute composante irréductible $X'_i$ de $X'$ de}}
\deleted{$\dim_x X = k+1$ \hand{;} et c'est aussi nécessaire si chaque $X'_i$ est non réglé rel.}
\deleted{à $f$} \inserted{\unsure{de dim $\leq k+1$ est, \deleted{\unsure{soit}} réglée rel. à $f$, \deleted{\unsure{soit}} telle que $X'_i = X_i$.}}
\inserted{\unsure{soit contenue} \unsure{dans $X'$}}

En effet, pour tout $\hand{$\zeta$}$, l'ensemble géométriquement singulier

de $Y_{\hand{$\zeta$}}$ \xout{xxxxxxxxxxxxxx(pt} (NB on se borne aux $\hand{$\zeta$}$ tels que $\phi_{\hand{$\zeta$}}$ soit $X_{k(\hand{$\zeta$})}$-

régulier, ce qui est inoffensif vu 8.2.) est la réunion de

sing$(Y'_{\hand{$\zeta$}}$ ) et de l'image inverse $T_{\hand{$\zeta$}}$ de $T$ dans $Y_{\hand{$\zeta$}}$ , donc la codimension
\inserted[typed]{l'inf de}
\margin{\unsure{lieu non-ordin. géométrique}}
de cet ensemble singulier dans $Y_{\hand{$\zeta$}}$ est égale à \hand{(}codim (sing$(Y'_{\hand{$\zeta$}})$,$Y'_{\hand{$\zeta$}}$ )

et de codim $(T_{\hand{$\zeta$}},Y_{\hand{$\zeta$}})$. Bornons nous aux $\hand{$\zeta$}$ tels que \deleted{$X - f(\underline{Y}'^{\,\mathrm{sup\ sing}}$}\hand{)}
\margin{$*$} \inserted{sing$(Y'_\zeta)$ soit fini}
\hand{/}(ce qui est inoffensif, car \deleted{par hyp} cela conduit à se mettre dans le

complémentaire d'un ensemble de codimension $\geq 2$\hand{\unsure{$\zeta$}$\}$}), \over{l}{L}es points singulier

géométriques de $Y'_{\hand{$\zeta$}}$ sont \inserted{alors} isolés\over{,}{.} La conclusion résulte facilement de

là , et de 8.4. \deleted{A vrai dire, on peut donner un énoncé plus satis-%
faisant, en disant: il f et suffit que pour toute composante irréd
$X'_i$ de $X'$ qui est réglée rel à f, on ait dim $X'_i$ \xout{kxx} (resp. dim $X'_i$
k+1), et que l'ensemble singulier ..}

Combinant 8.13. \deleted{et} 8.5.\inserted[typed]{et 8.6.}\hand{/}, on trouve à la façon habituelle:

\uline{Corollaire} 8.14. On suppose $f:X\rightarrow P$ non ramifié et $X$ sans pts isolés.

\margin{\unsure{n}} a) Supposons $X$ séparable sur $k$. Pour que l'ensemble \deleted{\unsure{E}} des $\hand{$\zeta \in$}$ $\overset{\hand{$\vee$}}{P}$

tels que $\phi_{\hand{$\zeta$}}$ soit $X_{k(\hand{$\zeta$})}$ -régulier et $Y_{\hand{$\zeta$}}$ soit séparable ait un complé-%

mentaire de codimension $\geq 2$, il f et s qu\deleted{'il n'existe aucun point
géom. singulier \inserted{\deleted{\unsure{isolé}}} de $x$ tel que $\dim_x X = 1$, \deleted{et que toute} \inserted{ni de composante}
\inserted{irréductible $X'_i$ de $X'$ de dim 1 ``non réglée'' pour $f$, et que l'ex}}

\mspage{48}
